\section{Value of computation and stopping}
\label{sec:C09:value}
Inference budget allocation is a sequential decision problem. Let $b$ be the controller's information state and $V_0(b)=\max_{a\in\mathcal A(b)}\mathbb E[U(a)\mid b]$ the value of selecting now. Computation $k$ costs $c_k$ in utility units and returns evidence $Z$, producing state $b^{k,Z}$. Its one-step value is
\begin{equation}
 \operatorname{VOC}_k(b)=
 \mathbb E_Z[V_0(b^{k,Z})]-V_0(b)-c_k.
\end{equation}
Under calibrated Bayesian updating, unchanged external state, and retained access to all original actions, the unpriced information term is nonnegative by the expectation-of-maximum inequality. A learned judge or an irreversible external operation does not inherit these assumptions. Verification, retrieval, resampling, and delegation must be assessed by their effect on the selected artifact, including actual controller and worker costs.

Stopping when every one-step value is nonpositive is a myopic rule, not generally an optimal stopping theorem. In a constructed example, the correct binary answer is the exclusive-or of two independent fair bits. Either observation alone leaves optimal accuracy at $1/2$; both raise it to one. If each observation costs 0.1 in the same utility scale as correctness, each one-step net value is $-0.1$, while obtaining both has net value $0.3$. Complementary computations can therefore justify continuing even when neither is individually valuable.

For a finite decision model with remaining resource budget $B$, deterministic resource cost $r_k>0$, a known computation-transition law, and terminal selection utility, let $W$ be the best continuation value. The recursion is
\begin{align}
 V(b,B)&=\max\{V_0(b),W(b,B)\},\\
 W(b,B)&=\max_{k:r_k\leq B}
 \left[-c_k+\mathbb E_Z V(b^{k,Z},B-r_k)\right].
\end{align}
Set $W(b,B)=-\infty$ if no computation is feasible. These equations specify the ideal allocation problem; a deployed controller generally lacks its transition law and calibrated utilities. The 2026 meta-reasoning case in Section~\ref{sec:C09:meta2026} supplies a concrete harness and artifact-level observations, rather than solving this recursion exactly~\cite{rewrite_metareasoning2026}. Evaluate stopping against realized improvements and deadlines, including observation staleness induced by further deliberation. A model's statement that more thought would help is not itself an estimate of the displayed value.

Repeatedly checking the same candidate can have little value when checks share a failure mode. An independent executable test may be more informative than another paraphrased self-critique. Conversely, a cheap local consistency check can eliminate an obvious error before an expensive environment call. The controller's action space should therefore expose cost and evidence type, not just a generic request to think harder. A budget measured in model calls can be misleading because one call may replay a huge context while another evaluates a short expression. Report tokens, external costs, and wall-clock time alongside the call count.

One operational realization is a myopic recurrence over a candidate graph $G_n$ and remaining resource budget $B_n$. Decision epochs are measured after the controller's proposal and scoring costs have been charged. Let $\beta_v\geq0$ reserve resources for the terminal check. For each proposed computation $k\in\mathcal K(G_n)$, $\widehat\Delta_n(k)$ estimates selection-utility gain, $\widehat c_n(k)$ its cost in the same utility units, and $\widehat r_n(k)\geq0$ its total resource charge through the next decision epoch. Define
\begin{equation}
\begin{aligned}
 \mathcal K_n
 &=\{k\in\mathcal K(G_n):\widehat r_n(k)\leq B_n-\beta_v\},\\
 k_n&\in\argmax_{k\in\mathcal K_n}
       [\widehat\Delta_n(k)-\widehat c_n(k)].
\end{aligned}
\label{eq:C09:computation-selection}
\end{equation}
The choice $k_n$ is defined only for a nonempty finite candidate set; fix a tie rule. The continuation indicator and stopping index are
\begin{equation}
\begin{aligned}
 d_n&=
 \begin{cases}
 \mathbf1\{\widehat\Delta_n(k_n)>\widehat c_n(k_n)\},
   &B_n>\beta_v,\ \mathcal K_n\ne\varnothing,\\
 0,&B_n\leq\beta_v\ \vee\ \mathcal K_n=\varnothing,
 \end{cases}\\
 \tau&=\inf\{n\geq0:d_n=0\}.
\end{aligned}
\label{eq:C09:computation-stop}
\end{equation}
For $n<\tau$, let $a_n$ be the artifact actually returned by an execution operator $\mathcal E$ that applies the required isolation or current authorization. Its evidence record $z_n$ retains source provenance $\sigma_n$, observed execution and recording cost $\rho_n\geq0$, and the applicable configuration $\nu_n$. Let $\gamma_{n+1}\geq0$ charge preparation of the next decision epoch:
\begin{equation}
\begin{aligned}
 a_n&=\mathcal E(k_n),\\
 z_n&=(a_n,\sigma_n,\rho_n,\nu_n),\\
 G_{n+1}&=\Psi(G_n,z_n),\\
 B_{n+1}&=B_n-\rho_n-\gamma_{n+1}.
\end{aligned}
\label{eq:C09:computation-transition}
\end{equation}
Here $\Psi$ adds the observed artifact and its evidence links to the graph. For finite $\tau$, a declared acceptance map $\mathcal S$ selects an artifact or returns the abstention marker $\bot$. The bounded terminal checker $\mathcal V$ returns a verdict $v_\tau$ and supporting record $e_\tau$; its domain includes $\bot$, which produces an unresolved verdict and an abstention record:
\begin{equation}
\begin{aligned}
 a_\tau^*&=\mathcal S(G_\tau),\\
 (v_\tau,e_\tau)&=\mathcal V(a_\tau^*;\beta_v).
\end{aligned}
\label{eq:C09:terminal-verification}
\end{equation}
The gain estimate remains a hypothesis to validate, and $\mathcal V$ establishes only its specified property. Predicted resource feasibility does not guarantee $\rho_n+\gamma_{n+1}\leq B_n-\beta_v$; caps or justified bounds on both execution and controller work are needed to preserve the reserve. This formulation makes those admissibility obligations explicit without identifying the myopic controller with the ideal Bellman solution.

